\section{An exact extension to taller symplectic corners}
\label{sec:taller}

The determinant in the two-plane theorem is a special case of a
symplectic area. This observation extends the exact law to every
proper symplectic target summand, when the input still consists of
one symplectic pair.

\begin{corollary}[A symplectic pair projected to any proper symplectic summand]
\label{cor:taller}
Let $1\le a<g$. Suppose $X^{\mathsf T}JX=J_2$ and
$Y^{\mathsf T}JY=J_{2a}$, where $Y$ has $2a$ columns and $J_{2a}$
is a standard unimodular alternating matrix. For
$M=Y^{\mathsf T}JSX\in\Mat_{2a,2}(\OO)$, put
\[
 \Omega(M)=M_1^{\mathsf T}J_{2a}M_2,\quad
 N=2g-2a,\quad
 \alpha_{g,a}=\frac{1-q^{-N}}{1-q^{-2g}}.
\]
The finite-precision density relative to $2a\times2$ independent entries is
\[
 h_{g,a,e}(M)=\frac{f_{N,e}(1-\Omega(M))}{1-q^{-2g}},
\]
with $f_{N,e}$ as in \eqref{eq:padic-finite-f}; the complete-ring
density is the corresponding geometric sum
\[
 \alpha_{g,a}\sum_{r=0}^{v(1-\Omega(M))}q^{-(N-1)r}.
\]
For every finite or infinite precision,
\begin{equation}\label{eq:taller-TV}
 \TV(\mu_e,\nu_e)
 =\left(1-q^{-1}+q^{-(2a+1)}\right)
   \frac{q^{2a}-1}{q^{2g}-1}.
\end{equation}
The event $\Omega(M)\equiv1\pmod\pi$ is optimal. For $e>k\ge1$,
\begin{equation}\label{eq:taller-lift}
 \TV(\mu_e,\widetilde\mu_{k,e})
 =(1-q^{-1})(1-\alpha_{g,a})q^{-Nk}.
\end{equation}
\end{corollary}

\begin{proof}
Split the ambient module as the target symplectic summand and its
symplectic complement. Write the image frame as $(B;C)$. Since
$M=J_{2a}B$, the visible pairing is
$\Omega(M)=B_1^{\mathsf T}J_{2a}B_2$, and the complement equation is
$\langle C_1,C_2\rangle=1-\Omega(M)$. The frame-counting proof of
Theorem~\ref{thm:padic-density} applies verbatim, with $N=2g-2a$;
the independent-entry factor is now $q^{-4ae}$.

For independent visible columns modulo $\pi$, the first is nonzero
with probability $1-q^{-2a}$, and then pairing with the second is
uniform. Consequently the area-one event has probability
$q^{-1}(1-q^{-2a})$. On its complement the likelihood is
$\alpha_{g,a}<1$; on the event it exceeds one by the same sign
calculation as before. Here $N\ge2$, which is the necessary inequality
in that calculation. This proves \eqref{eq:taller-TV}.

Conditional on an area-one residue class, one coordinate of the first
column is a unit. Varying the paired coordinate of the second column
makes the area uniform in $1+\pi R_e$. The valuation probabilities
used in Theorem~\ref{thm:padic-lift-exact} therefore hold with
$\beta_q=q^{-1}(1-q^{-2a})$. The identical calculation proves
\eqref{eq:taller-lift}.
\end{proof}

When $a=g$, there is no complementary symplectic pair and the
conclusion changes: the area is identically one, as in the $g=1$
endpoint of the square case. Thus the proper-summand hypothesis
has a concrete geometric role.
